% Осень26!B23; нумерация внутри столбца без пустых строк.
% Источник: лекция 4, §8 (строка 252)
\definitionticket{20}{Cтупенчатая функция, допустимое разбиение}
\label{definition:20}

\begin{definition}
\textbf{Разбиение} множества $E$ --- представление
\[
  E=\bigsqcup_{i=1}^{n}e_i
\]
в виде конечного объединения попарно непересекающихся подмножеств.
\end{definition}

\begin{definition}
Функция $f:E\to\mathbb R$ называется \textbf{ступенчатой}, если существует разбиение $E=\bigsqcup_{i=1}^{n}e_i$, на каждой части которого $f$ постоянна:
\[
  f|_{e_i}\equiv c_i,\qquad i=1,\ldots,n.
\]
Такое разбиение называется \textbf{допустимым} для $f$.
\end{definition}
